% GENERATED FILE. Edit canonical lessons, then run npm run generate:books.
% canonical-source-hash: fnv1a64:03e536014555d0a0
% canonical-lessons: TE-C405-arthika, TE-C405-dravyolbanam, TE-C405-pettubadi, TE-C405-badjettu, TE-C405-runam, TE-R405-review-398-405

\chapter{Money and the Economy}
\label{ch:te-money-and-the-economy}
\hlchaptermodality{405}

\begin{chapteropening}
\textbf{By the end of this chapter:} \emph{I can talk about the economy, prices, budgets and loans.}

Complete the last lesson of chapter 405: I can talk about the economy, prices, budgets and loans.
\end{chapteropening}

\section[\texorpdfstring{\te{ఆర్థిక} (ārthika)}{ఆర్థిక (ārthika)}]{\te{ఆర్థిక} (ārthika) — economic; financial}

\label{lesson:TE-C405-arthika}



\begin{quote}
Before the new one: say the Telugu for independent, then the Telugu for unemployment.
\end{quote}

\subsection*{You'll want to know: \te{ఆర్థిక}}
\textbf{\te{ఆర్థిక}} — \emph{ārthika} — \textquotedblleft{}economic; financial\textquotedblright{}.

\begin{quote}
\te{దేశం} \te{ఆర్థిక} \te{సమస్యల్లో} \te{ఉంది}.
\end{quote}

\begin{quote}
\emph{dēśaṁ ārthika samasyallō undi.}
\end{quote}

\begin{quote}
The country is in economic trouble.
\end{quote}

Before a noun: \te{ఆర్థిక} \te{సహాయం}, financial help. \te{ఆర్థికంగా}, financially.

\begin{cousinweb}
\textbf{\te{ఆర్థిక}} (\emph{ārthika}) is borrowed from Sanskrit \emph{ārthika}, \textquotedblleft{}significant; wise; rich\textquotedblright{} (Monier-Williams), from \emph{artha}, \textquotedblleft{}aim, wealth\textquotedblright{}; Telugu uses it for money matters.
\end{cousinweb}

\begin{grammarlens}[title={What you've built}]
One more word for giving reasons and opinions.
\end{grammarlens}

\subsection*{Your turn}
\begin{itemize}
\raggedright
  \item \emph{Say it:} \emph{ārthika}
  \item \emph{Say it:} \emph{ārthika}, once more
  \item \emph{Say it:} \emph{dēśaṁ ārthika samasyallō undi.} — The country is in economic trouble.
  \item \emph{Recall:} say the Telugu for independent, then the Telugu for unemployment, then say \emph{ārthika} again
\end{itemize}

\begin{tcolorbox}[breakable,colback=teal!4,colframe=teal!35!black,title={Before you move on}]
What does \te{ఆర్థిక} mean? (\textquotedblleft{}Economic; financial\textquotedblright{}.) Say it once more.
\end{tcolorbox}
\section[\texorpdfstring{\te{ద్రవ్యోల్బణం} (dravyōlbaṇaṁ)}{ద్రవ్యోల్బణం (dravyōlbaṇaṁ)}]{\te{ద్రవ్యోల్బణం} (dravyōlbaṇaṁ) — inflation (of prices)}

\label{lesson:TE-C405-dravyolbanam}



\begin{quote}
Before the new one: say the Telugu for unemployment, then the Telugu for economic.
\end{quote}

\subsection*{You'll want to know: \te{ద్రవ్యోల్బణం}}
\textbf{\te{ద్రవ్యోల్బణం}} — \emph{dravyōlbaṇaṁ} — \textquotedblleft{}inflation (of prices)\textquotedblright{}.

\begin{quote}
\te{ద్రవ్యోల్బణం} \te{వల్ల} \te{ధరలు} \te{పెరుగుతున్నాయి}.
\end{quote}

\begin{quote}
\emph{dravyōlbaṇaṁ valla dharalu perugutunnāyi.}
\end{quote}

\begin{quote}
Prices are rising because of inflation.
\end{quote}

A newspaper word; in talk, \te{ధరలు} \te{పెరగడం}, prices going up.

\begin{cousinweb}
\textbf{\te{ద్రవ్యోల్బణం}} (\emph{dravyōlbaṇaṁ}) is built from Sanskrit \emph{dravya}, \textquotedblleft{}substance, property\textquotedblright{} (CDIAL 6625), and \emph{ulbaṇa}, \textquotedblleft{}laid over in addition\textquotedblright{} (Monier-Williams); the compound is a modern one.
\end{cousinweb}

\begin{grammarlens}[title={What you've built}]
One more word for giving reasons and opinions.
\end{grammarlens}

\subsection*{Your turn}
\begin{itemize}
\raggedright
  \item \emph{Say it:} \emph{dravyōlbaṇaṁ}
  \item \emph{Say it:} \emph{dravyōlbaṇaṁ}, once more
  \item \emph{Say it:} \emph{dravyōlbaṇaṁ valla dharalu perugutunnāyi.} — Prices are rising because of inflation.
  \item \emph{Recall:} say the Telugu for unemployment, then the Telugu for economic, then say \emph{dravyōlbaṇaṁ} again
\end{itemize}

\begin{tcolorbox}[breakable,colback=teal!4,colframe=teal!35!black,title={Before you move on}]
What does \te{ద్రవ్యోల్బణం} mean? (\textquotedblleft{}Inflation (of prices)\textquotedblright{}.) Say it once more.
\end{tcolorbox}
\section[\texorpdfstring{\te{పెట్టుబడి} (peṭṭubaḍi)}{పెట్టుబడి (peṭṭubaḍi)}]{\te{పెట్టుబడి} (peṭṭubaḍi) — an investment; capital}

\label{lesson:TE-C405-pettubadi}



\begin{quote}
Before the new one: say the Telugu for economic, then the Telugu for inflation.
\end{quote}

\subsection*{You'll want to know: \te{పెట్టుబడి}}
\textbf{\te{పెట్టుబడి}} — \emph{peṭṭubaḍi} — \textquotedblleft{}an investment; capital\textquotedblright{}.

\begin{quote}
\te{అతను} \te{కొత్త} \te{దుకాణంలో} \te{పెట్టుబడి} \te{పెట్టాడు}.
\end{quote}

\begin{quote}
\emph{atanu kotta dukāṇaṁlō peṭṭubaḍi peṭṭāḍu.}
\end{quote}

\begin{quote}
He invested in a new shop.
\end{quote}

\te{పెట్టుబడి} \te{పెట్టు}, to invest; the first part is \te{పెట్టు}, to put.

\begin{cousinweb}
\textbf{\te{పెట్టుబడి}} (\emph{peṭṭubaḍi}) is native Telugu, built on \emph{peṭṭu}, \textquotedblleft{}to put, to place\textquotedblright{}; the origin of the second part, -\te{బడి}, is uncertain.
\end{cousinweb}

\begin{grammarlens}[title={What you've built}]
One more word for giving reasons and opinions.
\end{grammarlens}

\subsection*{Your turn}
\begin{itemize}
\raggedright
  \item \emph{Say it:} \emph{peṭṭubaḍi}
  \item \emph{Say it:} \emph{peṭṭubaḍi}, once more
  \item \emph{Say it:} \emph{atanu kotta dukāṇaṁlō peṭṭubaḍi peṭṭāḍu.} — He invested in a new shop.
  \item \emph{Recall:} say the Telugu for economic, then the Telugu for inflation, then say \emph{peṭṭubaḍi} again
\end{itemize}

\begin{tcolorbox}[breakable,colback=teal!4,colframe=teal!35!black,title={Before you move on}]
What does \te{పెట్టుబడి} mean? (\textquotedblleft{}An investment; capital\textquotedblright{}.) Say it once more.
\end{tcolorbox}
\section[\texorpdfstring{\te{బడ్జెట్టు} (baḍjeṭṭu)}{బడ్జెట్టు (baḍjeṭṭu)}]{\te{బడ్జెట్టు} (baḍjeṭṭu) — a budget}

\label{lesson:TE-C405-badjettu}



\begin{quote}
Before the new one: say the Telugu for inflation, then the Telugu for an investment.
\end{quote}

\subsection*{You'll want to know: \te{బడ్జెట్టు}}
\textbf{\te{బడ్జెట్టు}} — \emph{baḍjeṭṭu} — \textquotedblleft{}a budget\textquotedblright{}.

\begin{quote}
\te{మా} \te{ఇంటి} \te{బడ్జెట్టు} \te{చాలా} \te{చిన్నది}.
\end{quote}

\begin{quote}
\emph{mā iṇṭi baḍjeṭṭu cālā cinnadi.}
\end{quote}

\begin{quote}
Our household budget is very small.
\end{quote}

Like \te{బస్సు}, the English loan ends in -\te{ఉ}.

\begin{cousinweb}
\textbf{\te{బడ్జెట్టు}} (\emph{baḍjeṭṭu}) is a loan from English \emph{budget}, with the final -\te{ఉ} that many English loans take.
\end{cousinweb}

\begin{grammarlens}[title={What you've built}]
One more word for giving reasons and opinions.
\end{grammarlens}

\subsection*{Your turn}
\begin{itemize}
\raggedright
  \item \emph{Say it:} \emph{baḍjeṭṭu}
  \item \emph{Say it:} \emph{baḍjeṭṭu}, once more
  \item \emph{Say it:} \emph{mā iṇṭi baḍjeṭṭu cālā cinnadi.} — Our household budget is very small.
  \item \emph{Recall:} say the Telugu for inflation, then the Telugu for an investment, then say \emph{baḍjeṭṭu} again
\end{itemize}

\begin{tcolorbox}[breakable,colback=teal!4,colframe=teal!35!black,title={Before you move on}]
What does \te{బడ్జెట్టు} mean? (\textquotedblleft{}A budget\textquotedblright{}.) Say it once more.
\end{tcolorbox}
\section[\texorpdfstring{\te{రుణం} (ruṇaṁ)}{రుణం (ruṇaṁ)}]{\te{రుణం} (ruṇaṁ) — a loan; a debt}

\label{lesson:TE-C405-runam}



\begin{quote}
Before the new one: say the Telugu for an investment, then the Telugu for a budget.
\end{quote}

\subsection*{You'll want to know: \te{రుణం}}
\textbf{\te{రుణం}} — \emph{ruṇaṁ} — \textquotedblleft{}a loan; a debt\textquotedblright{}.

\begin{quote}
\te{రైతులు} \te{బ్యాంకు} \te{నుంచి} \te{రుణం} \te{తీసుకున్నారు}.
\end{quote}

\begin{quote}
\emph{raitulu byāṅku nuñci ruṇaṁ tīsukunnāru.}
\end{quote}

\begin{quote}
The farmers took a loan from the bank.
\end{quote}

Plural: \te{రుణాలు}, \emph{ruṇālu}.

Formal beside \te{అప్పు}: banks and governments give \te{రుణాలు}.

\begin{cousinweb}
\textbf{\te{రుణం}} (\emph{ruṇaṁ}) is borrowed from Sanskrit \emph{ṛṇa}, \textquotedblleft{}anything due, a debt, money owed\textquotedblright{} (Monier-Williams).
\end{cousinweb}

\begin{grammarlens}[title={What you've built}]
One more word for giving reasons and opinions.
\end{grammarlens}

\subsection*{Your turn}
\begin{itemize}
\raggedright
  \item \emph{Say it:} \emph{ruṇaṁ}
  \item \emph{Say it:} \emph{ruṇaṁ}, once more
  \item \emph{Say it:} \emph{raitulu byāṅku nuñci ruṇaṁ tīsukunnāru.} — The farmers took a loan from the bank.
  \item \emph{Recall:} say the Telugu for an investment, then the Telugu for a budget, then say \emph{ruṇaṁ} again
\end{itemize}

\begin{tcolorbox}[breakable,colback=teal!4,colframe=teal!35!black,title={Before you move on}]
What does \te{రుణం} mean? (\textquotedblleft{}A loan; a debt\textquotedblright{}.) Say it once more.
\end{tcolorbox}
\section[(dialogue)]{Words from chapters 398–405 — a review}

\label{lesson:TE-R405-review-398-405}



\begin{quote}
Say the Telugu for a law and to violate.
\end{quote}

\subsection*{Your turn}
Say these aloud:

\begin{itemize}
\raggedright
  \item a law, to violate, to implement, a punishment, to ban
  \item a duty, freedom, equality, injustice, to honour
  \item an election, a candidate, to elect, a campaign, a majority
  \item a citizen, democracy, a constitution, national, politics
  \item a legislative assembly, a minister, power, to govern, to approve
  \item administration, a representative, coordination, a scheme, welfare
  \item poverty, rich, wealth, independent, unemployment
  \item economic, inflation, an investment, a budget, a loan
\end{itemize}

\begin{tcolorbox}[breakable,colback=teal!4,colframe=teal!35!black,title={Before you move on}]
A law? (\textbf{\te{చట్టం}}, \emph{caṭṭaṁ}.) To violate? (\textbf{\te{ఉల్లంఘించు}}, \emph{ullaṅghiñcu}.) To implement? (\textbf{\te{అమలు} \te{చేయు}}, \emph{amalu cēyu}.) A punishment? (\textbf{\te{శిక్ష}}, \emph{śikṣa}.) To ban? (\textbf{\te{నిషేధించు}}, \emph{niṣēdhiñcu}.) A duty? (\textbf{\te{కర్తవ్యం}}, \emph{kartavyaṁ}.) Freedom? (\textbf{\te{స్వేచ్ఛ}}, \emph{svēccha}.) Equality? (\textbf{\te{సమానత్వం}}, \emph{samānatvaṁ}.) Injustice? (\textbf{\te{అన్యాయం}}, \emph{anyāyaṁ}.) To honour? (\textbf{\te{గౌరవించు}}, \emph{gauraviñcu}.) An election? (\textbf{\te{ఎన్నిక}}, \emph{ennika}.) A candidate? (\textbf{\te{అభ్యర్థి}}, \emph{abhyarthi}.) To elect? (\textbf{\te{ఎన్నుకో}}, \emph{ennukō}.) A campaign? (\textbf{\te{ప్రచారం}}, \emph{pracāraṁ}.) A majority? (\textbf{\te{మెజారిటీ}}, \emph{mejāriṭī}.) A citizen? (\textbf{\te{పౌరుడు}}, \emph{pauruḍu}.) Democracy? (\textbf{\te{ప్రజాస్వామ్యం}}, \emph{prajāsvāmyaṁ}.) A constitution? (\textbf{\te{రాజ్యాంగం}}, \emph{rājyāṅgaṁ}.) National? (\textbf{\te{జాతీయ}}, \emph{jātīya}.) Politics? (\textbf{\te{రాజకీయాలు}}, \emph{rājakīyālu}.) A legislative assembly? (\textbf{\te{శాసనసభ}}, \emph{śāsanasabha}.) A minister? (\textbf{\te{మంత్రి}}, \emph{mantri}.) Power? (\textbf{\te{అధికారం}}, \emph{adhikāraṁ}.) To govern? (\textbf{\te{పాలించు}}, \emph{pāliñcu}.) To approve? (\textbf{\te{ఆమోదించు}}, \emph{āmōdiñcu}.) Administration? (\textbf{\te{పరిపాలన}}, \emph{paripālana}.) A representative? (\textbf{\te{ప్రతినిధి}}, \emph{pratinidhi}.) Coordination? (\textbf{\te{సమన్వయం}}, \emph{samanvayaṁ}.) A scheme? (\textbf{\te{పథకం}}, \emph{pathakaṁ}.) Welfare? (\textbf{\te{సంక్షేమం}}, \emph{saṅkṣēmaṁ}.) Poverty? (\textbf{\te{పేదరికం}}, \emph{pēdarikaṁ}.) Rich? (\textbf{\te{ధనిక}}, \emph{dhanika}.) Wealth? (\textbf{\te{సంపద}}, \emph{sampada}.) Independent? (\textbf{\te{స్వతంత్ర}}, \emph{svatantra}.) Unemployment? (\textbf{\te{నిరుద్యోగం}}, \emph{nirudyōgaṁ}.) Economic? (\textbf{\te{ఆర్థిక}}, \emph{ārthika}.) Inflation? (\textbf{\te{ద్రవ్యోల్బణం}}, \emph{dravyōlbaṇaṁ}.) An investment? (\textbf{\te{పెట్టుబడి}}, \emph{peṭṭubaḍi}.) A budget? (\textbf{\te{బడ్జెట్టు}}, \emph{baḍjeṭṭu}.) A loan? (\textbf{\te{రుణం}}, \emph{ruṇaṁ}.)
\end{tcolorbox}
